\documentclass[11pt,letterpaper]{article}

\usepackage[T1]{fontenc}
\usepackage{amsmath,amssymb}
\usepackage{tikz}
\usepackage{pgfplots}
\pgfplotsset{compat=1.18}
\usepackage{algorithm}
\usepackage{algpseudocode}
\usepackage{booktabs}
\usepackage{enumitem}
\usepackage[hidelinks]{hyperref}

\newcommand{\N}{\mathbb{N}}
\newcommand{\Z}{\mathbb{Z}}
\newcommand{\Q}{\mathbb{Q}}
\newcommand{\R}{\mathbb{R}}
\newcommand{\C}{\mathbb{C}}
\newcommand{\Map}[3]{#1 \colon #2 \to #3}
\newcommand{\Mod}[1]{#1\text{-}\mathrm{Mod}}
\newcommand{\GroupHom}{\mathrm{Hom}_{\mathbf{Grp}}}
\newcommand{\RingHom}{\mathrm{Hom}_{\mathbf{Ring}}}

\newtheorem{theorem}{Theorem}
\newtheorem{lemma}[theorem]{Lemma}
\newtheorem{corollary}[theorem]{Corollary}
\newtheorem{definition}[theorem]{Definition}
\newtheorem{remark}[theorem]{Remark}
\newenvironment{proof}{\noindent\textit{Proof.}}{\hfill\(\square\)}

\title{Generalized Collatz Library}
\author{}
\date{\today}

\begin{document}
\maketitle

\begin{abstract}
For the positive-divisor, division-first raw map
\(T(n)=n/D\) when \(D\mid n\), and \(T(n)=qn+c\) otherwise,
assume \(q\ge2\), \(c\ge1\), \(D\ge2\), and \(\gcd(q,D)=1\).
Let \(m\) be the least positive return time of zero under the affine
residue permutation \(r\mapsto qr+c\bmod D\).
We prove that seeds with at least \(h\) actual raw division events
occupy exactly \(m^h\) classes modulo \(D^h\).
If \(m<D\), the seeds with infinitely many divisions have ordinary natural
density zero; the complementary density-one set has an exact affine tail
and escapes in the full raw clock.
At \(m=D\), every seed divides infinitely often.
The exceptional set need not be empty, and the return order need not divide \(D\).
All guard, lifting, raw-clock, density, and boundary statements have matching
Lean proofs. A separate exact Go census checks finite instances.
\end{abstract}

\section{Scope and provenance}

This is the positive scalar-divisor case of the residue-leakage claim in
the read-only historical corpus, Theorem 9.34 of
\texttt{FINALEE.MD}, lines 2019--2023 \cite{Corpus26}.
No negative divisor is replaced by its absolute value.
The original library admits signed parameters and signed states; here we
take the invariant class of positive states with \(D>0\).
Casting the natural-number states to integers commutes with every raw
step and every raw iterate. This is proved in
\texttt{RawCast} and \texttt{IterateCast}; it is not inferred from residue agreement.

The proof does not use an iid guard model, a distribution on one orbit,
a logarithmic average, or countable additivity of natural density.
Natural density here is the limit of initial-interval counts.
The finite computation below is supplementary evidence, not the
finite-to-infinite argument.

For \(j\ge0\), write
\[
G_j(q)=\sum_{i=0}^{j-1}q^i,\qquad G_0(q)=0.
\]
Define \(F(r)=qr+c\bmod D\) on \(\Z/D\Z\).
The prior residue family proves
\(F^j(0)\equiv cG_j(q)\pmod D\), the least-return criterion,
distinctness below the first return, and \(m\le D\).
The new proof \texttt{AccessExists} supplies existence for every unit \(q\bmod D\):
\(F\) is a permutation of a finite set, so zero is periodic.
Taking its least positive return gives a cycle of length \(1\le m\le D\).
In particular, an additional return-existence hypothesis is not being imposed.

\begin{definition}
The accessible set is
\[
A=\{F^j(0):0\le j<m\}.
\]
For \(h\ge0\), let \(S_h\) be the set of nonnegative integer seeds
whose raw orbit has at least \(h\) division events.
An event at raw time \(t\) is a division when \(D\mid T^t(n)\);
it includes an immediate division and consecutive divisions.
Let \(I=\bigcap_{h\ge0}S_h\).
\end{definition}

\section{A complete, bounded first-division search}

\begin{lemma}\label{LemmaSearch}
Starting from \(n\), the first division exists if and only if
\(n\bmod D\in A\). When it exists, it follows \(k\) affine events
for a unique first delay \(0\le k<m\), and the endpoint is
\[
H(n)=\frac{q^kn+cG_k(q)}{D}.
\]
All earlier affine states are not divisible by \(D\).
When the search through delays \(0,\ldots,m-1\) fails,
the same raw orbit never divides.
\end{lemma}
\begin{proof}
The affine state after \(k\) events is \(q^kn+cG_k(q)\), with residue
\(F^k(n\bmod D)\).
Since \(F\) is bijective and \(A\) is the cycle of zero,
a residue can reach zero if and only if it belongs to that cycle.
If \(n\bmod D=F^j(0)\), its first visit to zero occurs within one cycle.
This includes \(j=0\), for which the raw division is immediate.
Conversely, if \(F^k(n\bmod D)=0\), applying the inverse permutation \(k\) times
shows that \(n\bmod D\) belongs to the same cycle.
Thus a successful first visit has \(k<m\).
Choosing the first such \(k\) ensures every earlier raw step is the affine branch,
so the displayed endpoint is an integer and the raw time consumed is \(k+1\).
If no such \(k<m\) exists, no later affine visit to zero exists either.
Every subsequent raw event is then affine.
These are \texttt{HitAccessible}, \texttt{EventualAccessible},
\texttt{HitEventual}, \texttt{HitRawWitness}, and
\texttt{HitNoneNoDivisions}.
\end{proof}

\section{Jointly legal cylinders and source-digit lifting}

\begin{lemma}\label{LemmaShift}
Suppose \(h\) legal first-division blocks from \(n\) end at \(E\),
with numerator multiplier \(P=q^{k_1+\cdots+k_h}\).
For every \(u\ge0\), the source \(n+D^hu\) has precisely the same
first \(h\) block delays and endpoint \(E+Pu\).
Its multiplier is still \(P\), and \(\gcd(P,D)=1\).
Success or failure through depth \(h\) depends only on \(n\bmod D^h\).
\end{lemma}
\begin{proof}
For one block, adding \(Du\) preserves each division guard because
\[
q^k(n+Du)+cG_k(q)=q^kn+cG_k(q)+Dq^ku.
\]
The first divisible affine state therefore has the same delay.
Exact division changes the endpoint by \(q^ku\).
For \(h\) blocks, start with the shift \(D^hu\), apply the one-block identity,
and induct through the subsequent blocks; one factor of \(D\) is removed
per legal division, and every affine multiplier is a power of \(q\).
The final shift is \(Pu\).
All factors of \(P\) are units modulo \(D\), proving coprimality.
The same induction handles unsuccessful searches as well.
Writing \(n=r+D^h u\), \(0\le r<D^h\), proves the last assertion.
These identities are proved in \texttt{HitShift}, \texttt{RunShift},
\texttt{RunCoprime}, and \texttt{SurvivesMod}.
\end{proof}

\begin{theorem}\label{TheoremCylinders}
For every \(h\ge0\), \(S_h\) consists of exactly \(m^h\) residue
classes modulo \(D^h\).
Each legal depth-\(h\) source class has exactly \(m\) legal source-digit
lifts modulo \(D^{h+1}\).
The cylinder test is equivalent to having at least \(h\) actual raw division events.
\end{theorem}
\begin{proof}
At depth zero, every seed is legal and there is one class modulo one.
Fix a legal canonical representative \(r<D^h\), with endpoint \(E\) and
unit multiplier \(P\).
Its possible children are \(r+D^ht\), \(0\le t<D\).
By Lemma~\ref{LemmaShift}, the endpoint before the next block is \(E+Pt\).
The map
\[
t\longmapsto E+Pt\bmod D
\]
is a bijection of the residue set because \(\gcd(P,D)=1\).
Exactly the \(m\) residues in \(A\) permit the next division,
so this one common parent has exactly \(m\) legal children.
Children of distinct parents are disjoint: reducing a child modulo \(D^h\)
recovers its unique parent. Every legal depth-\(h+1\) source has this parent
and its unique high digit, so no cases are omitted.
The number of legal classes therefore satisfies
\(a_0=1\) and \(a_{h+1}=ma_h\), giving \(a_h=m^h\).
The Lean proof is \texttt{ChildrenCard}, \texttt{ChildrenDisjoint},
\texttt{ClassesSucc}, and \texttt{ClassesCard}.

Each successful block consumes its first delay plus one raw step
and contains exactly one actual division.
Composition recovers the original raw endpoint with exactly \(h\)
divisions; these are \texttt{HitRawCount} and \texttt{RunRaw}.
Conversely, a finite raw orbit with at least \(h\) divisions has a first
division whenever \(h>0\).
Lemma~\ref{LemmaSearch} recognizes that division within the complete bounded search.
Remove that block and induct on the remaining required division count.
This proves \texttt{RawCountSurvives} and the equivalence
\texttt{SurvivesRawCount}.
A returned block delay satisfies \(k<m\), as proved in \texttt{HitDelayBound}.
Hence a legal depth-\(h\) prefix has a raw witness time \(h\le t\le mh\),
with exactly \(h\) divisions and the same endpoint; this is \texttt{RunRawBudget}.
Division counts cannot decrease when the raw clock is extended.
The already proved converse therefore gives the exact finite-window equivalence
\[
n\in S_h\quad\Longleftrightarrow\quad
\text{at least \(h\) divisions occur within the first \(mh\) raw events}.
\]
This is \texttt{SurvivesRawWindow}; it justifies the independent accepted/rejected
raw-clock tests without an empirical search cap.
Finally \texttt{InfiniteSurvives} identifies the same infinite raw predicate
with survival at all finite depths.
\end{proof}

\section{Interval counts and the natural-density limit}

For integers \(X\ge0\) and \(L\ge1\), put
\[
B(X,L)=
\begin{cases}
0,&X=0,\\
\lfloor(X-1)/L\rfloor+1,&X>0.
\end{cases}
\]
This is \(\lceil X/L\rceil\), with its empty-interval value retained.

\begin{theorem}\label{TheoremDensity}
For every \(h\ge0\) and \(X\ge0\),
\[
\#(I\cap[0,X))\le \#(S_h\cap[0,X))
 \le \min\{X,m^hB(X,D^h)\}.
\]
For \(X>0\),
\[
\frac{\#(I\cap[0,X))}{X}
 \le \left(\frac mD\right)^h+\frac{m^h}{X}.
\]
If \(m<D\), then \(I\) has ordinary natural density zero.
Its complement has ordinary natural density one.
\end{theorem}
\begin{proof}
Write each \(n<X\) uniquely as \(n=r+D^h a\),
where \(0\le r<D^h\).
When \(n\in S_h\), \(r\) is one of the \(m^h\) legal canonical representatives.
Also \(0\le a<B(X,D^h)\): for \(X>0\), \(n\le X-1\);
for \(X=0\) there are no seeds to pack.
The injection \(n\mapsto(r,a)\) gives the ceiling estimate,
and the interval itself gives the bound by \(X\).
The inclusion \(I\subseteq S_h\) gives the first inequality.
These are \texttt{PrefixBelowCeilingBound} and
\texttt{InfiniteBelowCeilingBound}.
Since \(B(X,D^h)\le X/D^h+1\), division by \(X>0\)
gives the displayed boundary-term estimate
\texttt{InfiniteDensityBound}.

Fix the parameters and \(\varepsilon>0\).
First choose a finite depth \(h\) so that \((m/D)^h<\varepsilon/2\).
This is possible because \(0<m/D<1\).
Then choose an integer \(N>2m^h/\varepsilon\), \(N>0\).
For every integer \(X\ge N\), the boundary term is below \(\varepsilon/2\).
Thus the density ratio is below \(\varepsilon\) on the entire remaining sequence.
This is \texttt{ResidueDensityZero}; no exchange of an infinite intersection
with a density operation is used.
The exact partition of every finite interval into \(I\) and its complement
gives \texttt{FiniteDensityOne}.
\end{proof}

\begin{corollary}
Whenever \(X\le D^h\), \(\#(I\cap[0,X))\le m^h\).
In particular, the certificate can be selected by integer power comparisons,
without numerical logarithms.
\end{corollary}
\begin{proof}
For \(X=0\) the count is zero.
For \(0<X\le D^h\), \(B(X,D^h)=1\).
Apply Theorem~\ref{TheoremDensity}.
This is \texttt{InfiniteBelowScaleBound}.
\end{proof}

\begin{corollary}\label{CorollaryPositive}
For every integer \(X\ge0\) and \(h\ge0\), the original positive interval obeys
\[
\#(I\cap[1,X])\le \min\{X,m^hB(X,D^h)\}.
\]
For \(X\ge1\), let \(\alpha=\log m/\log D\). Then
\[
\#(I\cap[1,X])\le mX^{\alpha}.
\]
When \(m<D\), the infinitely-dividing positive seeds have ordinary natural density
zero, and the finitely-dividing positive seeds have ordinary natural density one.
\end{corollary}
\begin{proof}
For a positive source \(n\), use its unique positive block representative
\[
r=(n-1)\bmod D^h+1,\qquad a=\lfloor(n-1)/D^h\rfloor,
\qquad n=r+D^ha.
\]
This retains the residue-zero right endpoint \(r=D^h\) instead of the excluded zero seed.
The finite guard test for \(r\) is equivalent to that for \(n\) by the exact source shift.
The map from \([1,D^h]\) to residues modulo \(D^h\) is bijective:
its sole zero-residue representative is \(D^h\), and all other representatives
are their residues themselves.
At most \(m^h\) positive representatives are legal.
For \(1\le n\le X\), \(a<B(X,D^h)\).
The injective packing \(n\mapsto(r,a)\) gives the ceiling estimate,
while \([1,X]\) itself gives the minimum by \(X\).
These are \texttt{PositivePrefixBelow\allowbreak{}ScaleBound},
\texttt{PositivePrefixBelow\allowbreak{}CeilingBound}, and
\texttt{PositiveInfiniteBelowMinBound}.
In particular, \(X\le D^h\) gives a positive survivor count at most \(m^h\),
including equality at the right boundary.

Choose the least nonnegative \(h\) such that \(X\le D^h\).
Existence follows from the elementary induction \(k+1\le D^k\) for \(D\ge2\),
proved in \texttt{PositiveScaleExists}.
If \(h=0\), then \(X=1\); the count is at most one and \(1\le m\).
Otherwise \(h=k+1\) and minimality gives \(D^k<X\le D^{k+1}\).
Because \(m\ge1\), \(\alpha\ge0\).
The real-power identity \((D^k)^{\alpha}=m^k\) follows by taking logarithms;
\(\log D>0\) and no zero denominator is introduced.
Consequently
\[
\#(I\cap[1,X])\le m^{k+1}=m\,(D^k)^{\alpha}\le mX^{\alpha}.
\]
This includes \(m=1\), for which \(\alpha=0\).
The identities and comparison are \texttt{ResiduePowerIdentity} and
\texttt{PositiveInfinitePowerBound}; the executable certificate uses only
integer scales, not rounded logarithms.

Finally \(I\cap[1,X]\subseteq I\cap[0,X+1)\).
For \(X\ge1\), the normalization factor \((X+1)/X\) is at most two.
Apply the already proved density-zero contract with tolerance \(\varepsilon/2\)
to \(X+1\). The result is below \(\varepsilon\) for every sufficiently large \(X\).
The positive finite and infinite counts sum to exactly \(X\), proving density one.
These are \texttt{PositiveResidueDensityZero}, \texttt{PositivePartition},
and \texttt{PositiveFiniteDensityOne}.
\end{proof}

\section{Recovery of every escaping raw tail}

\begin{theorem}\label{TheoremTail}
If \(n\notin I\), there is one raw time \(t_0\) after which no division occurs.
For every \(j\ge0\),
\[
T^{t_0+j}(n)=q^jT^{t_0}(n)+cG_j(q).
\]
For \(c\ge1\) and \(q\ge2\), this same full raw sequence tends to \(+\infty\).
Consequently the positive seeds that eventually stop dividing and then escape
have ordinary natural density one when \(m<D\).
Every eventually periodic positive seed belongs to \(I\).
\end{theorem}
\begin{proof}
The raw division count is finite for \(n\notin I\).
Equivalently, some complete finite-depth block search fails.
Recover the actual endpoint of every successful preceding block.
The failed block has no eventual division by Lemma~\ref{LemmaSearch},
so its endpoint begins one division-free raw tail.
This is \texttt{RunNoneTail} and \texttt{FiniteDivisionsTail}.
Every later raw event is affine. Induction gives the exact displayed formula,
using \(G_{j+1}(q)=qG_j(q)+1\); this is
\texttt{DivisionFreeAffine} and \texttt{FiniteDivisionsFormula}.

The affine state increases by at least one each step:
\(qx+c\ge x+1\) for \(x\ge0\).
Hence it eventually exceeds every fixed bound at every later raw clock,
as proved in \texttt{AffineGrowth}, \texttt{DivisionFreeEscape}, and
\texttt{ResidueEscape}.
For positive initial states, division by \(D\) when legal preserves positivity,
and the affine branch preserves positivity; see \texttt{IteratePositive}.
The complementary density-one conclusion follows from
Theorem~\ref{TheoremDensity}, not from an average orbit drift.
An eventually periodic orbit is bounded, whereas every seed outside \(I\) escapes.
Therefore every eventually periodic positive seed lies in \(I\).
The formal statements \texttt{BoundedInfinite} and \texttt{PeriodicInfinite}
retain the actual raw clock and a strictly positive period.
\end{proof}

\section{Sharp boundary and separating examples}

\begin{theorem}\label{TheoremBoundary}
If \(m=D\), every nonnegative seed belongs to \(I\).
This does not decide whether those infinitely-dividing orbits are bounded or escaping.
\end{theorem}
\begin{proof}
Theorem~\ref{TheoremCylinders} gives \(D^h\) legal classes modulo \(D^h\),
so the legal set is the entire residue range at every depth.
By the raw-event equivalence, every seed has arbitrarily many actual divisions.
These are \texttt{FullAccessClasses} and \texttt{FullAccessInfinite}.
\end{proof}

For \((q,c,D)=(2,1,3)\), the zero residue cycle is
\(0\mapsto1\mapsto0\); thus \(m=2<D=3\).
Yet the raw cycle \(1\mapsto3\mapsto1\) belongs to \(I\).
In fact the depth-\(h\) block from \(1\) returns to \(1\) with multiplier \(2^h\).
These are \texttt{TernaryPeriod}, \texttt{TernaryCycleRun},
and \texttt{TernaryCycleInfinite}.
A zero-density exceptional set can therefore be nonempty, and an atomic
source supported on this cycle assigns it full mass.
The survival fraction is not a capture probability.

The same example refutes the source's stronger assertion
\emph{order divides \(D\)}: \(m=2\) does not divide \(D=3\).
The valid cardinal bound is \(m\le D\).
The formal common-witness counterexample is
\texttt{PeriodNeedNotDivideModulus}.
For classical \((3,1,2)\), \(m=D=2\); see \texttt{ClassicalPeriod}.
The density-one finite-division theorem does not apply to that map,
and the full-access result does not prove classical convergence or divergence.

\section{Reproducible finite evidence and formal correspondence}

The source files are:
\begin{itemize}
\item \path{QCD.ResidueLeakage.001/lean/ResidueCylinders.lean};
\item \path{QCD.ResidueLeakage.001/lean/ResidueDensity.lean};
\item \path{QCD.ResidueLeakage.001/go/ResidueCylinders.go}.
\end{itemize}
The default Go census uses \(2\le q\le8\), \(2\le D\le9\),
\(0\le c\le3\), \(\gcd(q,D)=1\), and recorded depths \(0\le h\le4\).
Each record also checks all high-digit children at depth \(h+1\), including
independent raw-clock rejection checks over the complete \(mh\) or \(m(h+1)\) window.
The executed run produced 132 models, 660 records, 75,490 legal parent classes,
898,742 parent/child raw-clock checks, and 589,658 source-digit lift checks.
The zero-offset cases test the broader finite algebraic lemmas;
the positive-offset escape theorem retains \(c\ge1\).

Eight focused tests include immediate and consecutive divisions,
the nondividing residue \(2\bmod3\), strict and critical censuses,
a 256-bit shifted source, finite-interval boundaries, invalid parameters and
budgets, zero hot-loop allocations, and deterministic parallel assembly.
The finite report does not compute membership of an infinite predicate.
The universal recovery and limit statements are proved in Lean,
not extrapolated from this census.

The source coverage is limited to the positive-divisor density-one half
of Theorem 9.34, its finite cylinder and survivor estimates, and the explicit
return-order counterexample.
The conditional mean-slope and pressure statements in that source paragraph,
other source results, negative divisors, and the broader corpus-to-result
crosswalk are not declared complete by this family.

\bibliographystyle{amsalpha}
\bibliography{References}
\end{document}
